% Generated draft: requires mathematical and source review.
\documentclass[11pt]{article}
\usepackage[margin=25mm]{geometry}
\usepackage{fontspec}
\setmainfont{Latin Modern Roman}
\usepackage{amsmath,amssymb,mathtools,unicode-math}
\setmathfont{Latin Modern Math}
\usepackage{enumitem,needspace,xurl,hyperref}
\setlength{\parindent}{0pt}
\setlength{\parskip}{5pt}
\setlength{\emergencystretch}{3em}
\newcommand{\N}{\mathbb{N}}
\newcommand{\Z}{\mathbb{Z}}
\newcommand{\Q}{\mathbb{Q}}
\newcommand{\R}{\mathbb{R}}
\newcommand{\C}{\mathbb{C}}
\newcommand{\F}{\mathbb{F}}
\newcommand{\A}{\mathbb{A}}
\newcommand{\PP}{\mathbb{P}}
\newcommand{\E}{\mathbb{E}}
\newcommand{\eps}{\varepsilon}
\newcommand{\dd}{\,\mathrm{d}}
\begin{document}
\section*{UnsolvedMath ID: 30002335}
\textbf{Source problem code:} catalogue entry
\subsection{Definitions and mathematical statement}
A classical knot type \(K\) is an ambient-isotopy class of embeddings of the circle \(S^1\) into \(\R^3\). An open rope instead has a centreline parametrized by an interval. Associating such a centreline with a classical knot type requires an additional convention, such as a specified closure or restrictions on endpoint motions. The supplied record does not provide this convention.

For the purpose of specifying one possible angle observable, suppose a model represents a rope by a regular \(C^1\) curve \(\gamma:[0,L]\to\R^3\), parametrized by arclength, with \(L>0\). Its outward endpoint tangent directions are
\[
 u_0=-\gamma'(0),\qquad u_L=\gamma'(L).
\]
The angle between these directed vectors would be
\[
 \theta(\gamma)=\arccos\bigl(u_0\cdot u_L\bigr)\in[0,\pi].
\]
This is a conditional definition, not an angle convention or regularity hypothesis supplied by the source. The source does not specify whether its angle is measured between endpoint tangents, approximately straight free tails, unoriented lines, or other directions.

\textbf{Problem as supplied.} What mathematical model describes a tight open knot whose released ends form a characteristic angle, and how does that angle depend on the knot type?

\textbf{Conditional mathematical formulation.} Specify a model \(\mathcal M\) defining admissible open-rope configurations, their knot labels, the meaning of tightness, the release procedure, and the resulting angle observable. If \(p\) denotes the additional physical or geometric parameters and boundary data required by that model, let \(\mathcal R_{\mathcal M}(K;p)\) denote its resulting released configurations bearing label \(K\). Where an angle observable \(\theta_{\mathcal M}\) is defined, the possible angles form
\[
 \Theta_{\mathcal M}(K;p)
 =\{\theta_{\mathcal M}(\gamma):\gamma\in\mathcal R_{\mathcal M}(K;p)\}.
\]
The requested task is to identify an appropriate model and determine the dependence of its end-angle predictions on \(K\), including whether a single characteristic angle is well defined and which additional data affect it. No assertion that this set is nonempty or a singleton, or that its value depends only on \(K\), is supplied. Without choices of admissibility, tightness, release dynamics or equilibrium selection, endpoint conditions, and angle convention, the record does not determine a unique mathematical statement with fixed hypotheses and quantifiers.
\subsection{Short English statement}
The problem asks for a mathematical model predicting the angle made by the released ends of a tight open knot and explaining its dependence on knot type. A precise formulation still requires a convention for the knot type of an open rope, a definition of tightness, a release procedure, boundary conditions, and an angle measurement. The supplied record does not establish that knot type alone determines one angle.
\subsection{Sources}
\begin{itemize}
\item Local UnsolvedMath registry, \texttt{data/unsolvedmath/problems.json}, permanent ID \(30002335\), title \emph{End Angles of Tight Open Knots}, exact source code \texttt{catalogue entry}. The supplied record lists proposed year \(2013\) and gives no named proposer. Its statement status is recorded as \texttt{exact}; its verification field reports a manual assessment of self-containment, coherence, and absence of extraction artifacts. These are supplied curation claims, not independent mathematical verification.
\item Supplied report citation: \emph{Geometric Knot Theory (2014)}, Oberwolfach Reports, DOI \texttt{10.4171/owr/2013/22}, pp.~1313--1358. \url{https://doi.org/10.4171/owr/2013/22}. The same reference appears in the supplied literature sources as \emph{Oberwolfach Report 22/2013}. The differing year designations are retained without resolution.
\item Cantarella--Kusner--Sullivan, \emph{On the Minimum Ropelength of Knots and Links}. \url{https://arxiv.org/abs/math/0103224}.
\item Literature assessment dated exactly \texttt{2026-08-21}: the supplied curation labels the problem open and states that existence and \(C^{1,1}\) regularity of closed ropelength minimizers are known, while no general model or formula determines the end angle of a tight open knotted rope from its knot type. This is supplied dated curation, not independently verified. The background expressly requests independent verification of time-sensitive status and cited claims; the closed-rope statements do not supply a model for the open-rope question.
\end{itemize}
\end{document}
